% =============================================================================
% CKKS para Aprendizado Federado — guia acessível para praticantes de FL
% Público-alvo: pesquisadores de FL sem conhecimento prévio de HE.
% Compilar: pdflatex ckks_para_fl.tex  (2 passagens)
% =============================================================================
\documentclass[11pt,a4paper]{article}

\usepackage[utf8]{inputenc}
\usepackage[T1]{fontenc}
\usepackage[brazil]{babel}

\usepackage{amsmath}
\usepackage{amssymb}
\usepackage{amsthm}

\usepackage{tikz}
\usetikzlibrary{arrows.meta, positioning, calc, decorations.pathreplacing,
                shapes.geometric, fit, backgrounds}

\usepackage{booktabs}
\usepackage[a4paper, top=2.5cm, bottom=2.5cm, left=2.8cm, right=2.8cm]{geometry}

\usepackage[colorlinks=true, allcolors=blue,
            pdfauthor={Yuri Dimitre Dias de Faria},
            pdftitle={CKKS para Aprendizado Federado}]{hyperref}
\usepackage[brazilian, nameinlink]{cleveref}

\newtheorem{definition}{Definição}

% =============================================================================
\begin{document}

\title{\textbf{Criptografia Homomórfica CKKS\\
para Aprendizado Federado Preservador de Privacidade}\\[0.4em]
\large Guia conceitual para praticantes de FL}
\author{Yuri Dimitre Dias de Faria\\
\small Instituto de Computação -- UNICAMP\\
\small Orientador: Prof. Dr. Allan Mariano de Souza}
\date{Junho de 2026}
\maketitle

% =============================================================================
\section{O problema: privacidade na agregação federada}
\label{sec:problema}
% =============================================================================

No Aprendizado Federado (FL), um servidor central coordena o treinamento de um
modelo global sem coletar os dados brutos dos clientes. Cada cliente $k$ computa
um gradiente $g_k$ sobre seus dados locais e o transmite para o servidor, que
agrega as atualizações e devolve o novo modelo. O pressuposto de privacidade é que
apenas os gradientes, e não os dados de treinamento, saem dos dispositivos.

Esse pressuposto é fraco. Ataques de inversão de gradiente~\cite{Geiping2020}
reconstroem imagens de entrada com alta fidelidade a partir de $g_k$. Ataques de
inferência de pertença~\cite{Shokri2017} determinam se um exemplo específico
participou do treinamento. Em aplicações médicas, financeiras ou governamentais,
onde a lei exige proteção dos dados individuais, expor os gradientes ao servidor é
juridicamente problemático mesmo quando o servidor é considerado honesto.

A Criptografia Homomórfica (HE) elimina esse vetor de ataque: os clientes enviam
\emph{gradientes cifrados}, o servidor agrega matematicamente sem acessar os valores,
e apenas o detentor da chave secreta decifra o resultado final. A propriedade
fundamental que torna isso possível é o \textbf{homomorfismo}: a soma dos textos
cifrados, ao ser decifrada, produz a soma dos textos claros correspondentes,
sem que o servidor precise de qualquer chave.

% =============================================================================
\section{Criptografia Homomórfica: o essencial em uma página}
\label{sec:he_basico}
% =============================================================================

\paragraph{Texto claro, texto cifrado e par de chaves.}
Um esquema de criptografia homomórfica define três objetos fundamentais. O
\textbf{texto claro} (plaintext) é o dado original, por exemplo o vetor de
gradientes $g_k \in \mathbb{R}^d$. O \textbf{texto cifrado} (ciphertext) é uma
representação cifrada de $g_k$, legível apenas por quem detém a chave secreta. O
\textbf{par de chaves} consiste em uma chave pública $\text{pk}$, que qualquer
cliente usa para cifrar, e uma chave secreta $\text{sk}$, que apenas a entidade de
confiança mantém para decifrar. A assimetria é a garantia de privacidade: clientes
cifram com $\text{pk}$, o servidor manipula textos cifrados, e apenas o detentor de
$\text{sk}$ vê os valores.

\paragraph{O homomorfismo aditivo.}
A propriedade central para FL é a \textbf{adição homomórfica}: dado um operador de
adição $\oplus$ sobre textos cifrados,
\begin{equation}
  \mathrm{Dec}\!\left(\mathrm{Enc}(g_1) \oplus \mathrm{Enc}(g_2)\right)
  \approx g_1 + g_2,
  \label{eq:homo_add}
\end{equation}
onde $\mathrm{Enc}$ e $\mathrm{Dec}$ são as operações de cifragem e decifragem,
respectivamente. A igualdade em~\eqref{eq:homo_add} é aproximada: todo esquema de
HE prático incorpora um pequeno ruído criptográfico cujos efeitos são discutidos na
\cref{sec:custos}.

\paragraph{O esquema CKKS.}
O CKKS~\cite{Cheon2017}, também denominado HEAAN na literatura, é o esquema de HE
mais adequado para FL porque opera nativamente sobre \textbf{números reais e complexos
aproximados}. Outros esquemas (BFV, BGV) operam sobre inteiros exatos, o que obriga
a quantização prévia dos gradientes com perda de precisão adicional. O CKKS codifica
os gradientes diretamente com um fator de escala $\Delta$ e aceita um pequeno erro de
aproximação como custo intrínseco, similar ao erro de ponto flutuante de 64 bits.

\paragraph{Empacotamento de múltiplos valores por criptograma.}
O CKKS não cifra um único valor por vez. Cada criptograma acomoda simultaneamente
$N/2$ valores independentes, denominados \textbf{slots}, onde $N$ é o grau do
polinômio anel usado internamente. Para $N=4096$, um criptograma carrega 2048 gradientes;
para $N=8192$, carrega 4096. Isso reduz o número de operações de rede e de cifragem
na proporção $N/2$ em relação a cifrar escalar a escalar.

\paragraph{Por que só adição é suficiente para FL.}
O algoritmo FedAvg, e mais geralmente qualquer forma de agregação fedrerada por soma
ponderada, calcula
\[
  G = \sum_{k=1}^{K} w_k\, g_k,
\]
onde os pesos $w_k = n_k / \sum_j n_j$ são escalares conhecidos pelo servidor. Essa
operação é inteiramente aditiva do ponto de vista criptográfico: basta que o cliente
$k$ pré-multiplique $g_k$ por $w_k$ \emph{antes} de cifrar, e o servidor some os
criptogramas. Nenhuma multiplicação entre dois textos cifrados é necessária. Isso é
relevante porque multiplicação homomórfica entre dois textos cifrados é dispendiosa:
requer relinearização e reescalonamento (\emph{rescaling}), que aumentam o custo de
processamento e consomem o orçamento de ruído do esquema. A restrição ao regime
somente-aditivo elimina completamente essas operações.

% =============================================================================
\section{O fluxo do FL cifrado}
\label{sec:fluxo}
% =============================================================================

\paragraph{Preparação da rodada.}
Uma entidade de confiança (que pode ser o próprio servidor sob determinadas
premissas de segurança, ou uma terceira parte, ou um comitê de clientes) gera o par
de chaves $(\text{pk}, \text{sk})$. A chave pública $\text{pk}$ é distribuída a todos
os participantes; a chave secreta $\text{sk}$ permanece protegida.

\paragraph{Cifragem no cliente.}
Cada cliente $k$ computa o gradiente local $g_k \in \mathbb{R}^d$, divide-o em
segmentos de tamanho $N/2$, codifica cada segmento como um polinômio com escala
$\Delta$, e cifra com $\text{pk}$. O resultado é uma lista de textos cifrados
$\{\text{ct}_k^{(i)}\}$. O cliente envia apenas esses criptogramas ao servidor;
$g_k$ nunca sai em texto claro.

\paragraph{Agregação no servidor.}
O servidor recebe $\{\text{ct}_k^{(i)}\}$ de todos os $K$ clientes e soma os
criptogramas correspondentes ao mesmo segmento $i$:
\begin{equation}
  \text{ct}_{\text{agg}}^{(i)} = \bigoplus_{k=1}^{K} \text{ct}_k^{(i)}.
  \label{eq:agg}
\end{equation}
A operação $\oplus$ em~\eqref{eq:agg} é a adição coeficiente a coeficiente dos
polinômios que compõem cada criptograma, que não requer $\text{sk}$. A
\cref{fig:fl_pipeline} ilustra o fluxo completo.

\paragraph{Decifragem e atualização.}
O servidor envia $\{\text{ct}_{\text{agg}}^{(i)}\}$ ao detentor de $\text{sk}$, que
decifra, decodifica e recupera uma aproximação de $\sum_k g_k$. Essa agregação é
usada para atualizar o modelo global, que é distribuído em texto claro para a próxima
rodada.

\begin{figure}[ht]
\centering
\begin{tikzpicture}[
  node distance=0.4cm and 0.7cm,
  box/.style  = {draw, rounded corners=3pt, minimum width=2.1cm, minimum height=0.6cm,
                 align=center, font=\small},
  sbox/.style = {draw, rounded corners=3pt, minimum width=2.1cm, minimum height=0.6cm,
                 align=center, font=\small, fill=gray!10},
  io/.style   = {draw, ellipse, minimum width=1.5cm, minimum height=0.6cm,
                 align=center, font=\small},
  arr/.style  = {-{Stealth[length=5pt]}, thick},
  lab/.style  = {font=\scriptsize, text=gray, align=center}
]

% --- Cliente 1 ---
\node[io]  (g1)  {$g_1$};
\node[box, right=0.5cm of g1]  (c1)  {Codificar\\+ Cifrar($\text{pk}$)};
\node[io,  right=0.5cm of c1]  (ct1) {$\text{ct}_1$};

% --- Cliente 2 ---
\node[io, below=0.5cm of g1]  (g2)  {$g_2$};
\node[box, right=0.5cm of g2] (c2)  {Codificar\\+ Cifrar($\text{pk}$)};
\node[io,  right=0.5cm of c2] (ct2) {$\text{ct}_2$};

% --- Reticências ---
\node[below=0.35cm of g2, font=\small] (dots) {$\vdots$};
\node[below=0.5cm of g2] (gdots) {};

% --- Cliente K ---
\node[io, below=1.1cm of g2]  (gK)  {$g_K$};
\node[box, right=0.5cm of gK] (cK)  {Codificar\\+ Cifrar($\text{pk}$)};
\node[io,  right=0.5cm of cK] (ctK) {$\text{ct}_K$};

% --- Servidor ---
\node[sbox, right=1.2cm of ct2, yshift=-0.2cm] (srv) {Agregar\\$\oplus_{k=1}^K \text{ct}_k$};
\node[io,  right=0.55cm of srv]  (cta) {$\text{ct}_{\text{agg}}$};

% --- Decifrador ---
\node[box, right=0.55cm of cta] (dec) {Decifrar($\text{sk}$)\\+ Decodificar};
\node[io,  right=0.5cm of dec]  (G)   {$\sum_k g_k$};

% Flechas clientes
\draw[arr] (g1) -- (c1);
\draw[arr] (c1) -- (ct1);
\draw[arr] (ct1) -| (srv);

\draw[arr] (g2) -- (c2);
\draw[arr] (c2) -- (ct2);
\draw[arr] (ct2) -- (srv);

\draw[arr] (gK) -- (cK);
\draw[arr] (cK) -- (ctK);
\draw[arr] (ctK) -| (srv);

% Flechas servidor → decifrador
\draw[arr] (srv) -- (cta);
\draw[arr] (cta) -- (dec);
\draw[arr] (dec) -- (G);

% Labels
\node[lab, above=0.05cm of c1]   {cliente 1};
\node[lab, above=0.05cm of c2]   {cliente 2};
\node[lab, above=0.05cm of cK]   {cliente $K$};
\node[lab, above=0.05cm of srv]  {servidor\\(sem $\text{sk}$)};
\node[lab, above=0.05cm of dec]  {agregador\\(com $\text{sk}$)};
\end{tikzpicture}
\caption{Fluxo do FL com cifragem homomórfica. O servidor agrega criptogramas sem
acesso à chave secreta; o gradiente em texto claro nunca trafega pela rede.}
\label{fig:fl_pipeline}
\end{figure}

% =============================================================================
\section{O que custa: precisão, ruído e comunicação}
\label{sec:custos}
% =============================================================================

\paragraph{Escala $\Delta$ e erro de codificação.}
Antes de cifrar, cada valor real $g \in \mathbb{R}$ é multiplicado por $\Delta$ e
arredondado para um inteiro. O erro introduzido por esse arredondamento é delimitado
por
\begin{equation}
  |\varepsilon_{\text{cod}}| \leq \frac{1}{2\Delta}.
  \label{eq:enc_error}
\end{equation}
Aumentar $\Delta$ reduz o erro de codificação, mas exige um módulo aritmético $q$
maior para que os coeficientes do polinômio permaneçam dentro do intervalo válido.
Valores típicos são $\Delta = 2^{40}$ a $2^{50}$, garantindo de 12 a 15 dígitos
decimais de precisão antes da agregação.

\paragraph{Ruído criptográfico e crescimento com $K$.}
Todo esquema RLWE incorpora um ruído $e$ nos coeficientes do criptograma durante a
cifragem. Esse ruído é intencional: ele é a razão pela qual o adversário não consegue
recuperar a chave secreta. Na decifragem, o ruído criptográfico se soma ao erro de
codificação conforme a equação de erro total por slot:
\begin{equation}
  \hat{g} = g + \varepsilon_{\text{cod}} + \frac{e}{\Delta},
  \label{eq:total_error}
\end{equation}
onde $e$ é amostrado de uma distribuição gaussiana discreta com desvio padrão $\sigma_e$.

Ao agregar $K$ clientes, os ruídos individuais $e^{(k)}$ são independentes e de média
zero. A variância total é a soma das variâncias; portanto, o desvio padrão do erro
agregado por slot cresce como
\begin{equation}
  \sigma_{\text{agg}} \approx \frac{\sqrt{K}\,\sigma_e}{\Delta}.
  \label{eq:noise_agg}
\end{equation}
O crescimento sublinear de $\sqrt{K}$ é uma vantagem em relação a mecanismos de ruído
diferencial, que adicionam ruído proporcional a $K$. Ainda assim, em federações com
centenas de clientes, o termo $\sqrt{K}\,\sigma_e/\Delta$ deve permanecer muito abaixo
do sinal útil, o que impõe um limite inferior sobre $\Delta$.

\paragraph{Grau $N$: capacidade, segurança e custo de banda.}
O parâmetro $N$ regula três aspectos do sistema de forma simultânea. Primeiro, o
número de slots por criptograma é $N/2$: para um gradiente de dimensão $d$, o cliente
transmite $\lceil d/(N/2)\rceil$ criptogramas. Segundo, a resistência a ataques contra
o problema RLWE subjacente aumenta com $N$ para um módulo $q$ fixo; valores
comumente recomendados são $N \in \{4096, 8192, 16384\}$ para níveis de segurança de
128, 192 e 256 bits, respectivamente. Terceiro, o tamanho de cada criptograma é
$2 \times (12 + 8N)$ bytes, o que para $N=8192$ resulta em aproximadamente 131\,KB
por criptograma; a expansão de tamanho é o principal custo de comunicação da abordagem.

A \cref{tab:tradeoffs} resume os efeitos de variar cada parâmetro.

\begin{table}[ht]
\centering
\caption{Efeitos de aumentar cada parâmetro do CKKS no contexto de FL.}
\label{tab:tradeoffs}
\small
\begin{tabular}{lccc}
\toprule
\textbf{Parâmetro} & \textbf{Precisão} & \textbf{Segurança} & \textbf{Custo de banda} \\
\midrule
$\Delta$ maior   & melhora ($\uparrow$) & neutro          & neutro (precisa $q$ maior) \\
$N$ maior        & mais slots           & melhora ($\uparrow$) & aumenta ($\uparrow$) \\
$K$ maior        & piora ($\sqrt{K}$)   & neutro           & aumenta ($\uparrow$) \\
\bottomrule
\end{tabular}
\end{table}

\paragraph{Comparação com alternativas.}
A Privacidade Diferencial (DP) oferece proteção formal a custo computacional baixo, mas
injeta ruído que escala linearmente com $K$ e não oculta o valor exato do gradiente
do servidor. A Computação Segura Multipartidária (MPC) protege contra servidores
desonestos sem expansão de mensagem, mas exige múltiplas rodadas de comunicação. O
HE com CKKS é preferível quando o servidor deve ser incapaz de observar qualquer
gradiente individual, a latência de múltiplas rodadas é proibitiva, e a precisão em
ponto flutuante importa.

% =============================================================================
\section{Fundação de segurança e a biblioteca implementada}
\label{sec:seguranca}
% =============================================================================

\paragraph{O problema RLWE.}
A segurança do CKKS repousa sobre a hipótese do \textbf{Aprendizado com Erros em Anéis}
(RLWE)~\cite{LPR2010}. Dado um polinômio $a$ sorteado uniformemente e um polinômio de
ruído $e$ com coeficientes gaussianos pequenos, o RLWE afirma que o par
$(a,\, a \cdot s + e)$, onde $s$ é a chave secreta, é computacionalmente
indistinguível de dois polinômios uniformes independentes. Em outras palavras: mesmo
observando o criptograma $(c_0, c_1)$ inteiro, um adversário sem $\text{sk}$ não
consegue recuperar $g_k$ em tempo polinomial, assumindo que RLWE é intratável. A
escolha dos parâmetros $(N, q)$ é guiada por estimativas de custo de ataques de
retículo (LWE Estimator), e valores como $N=8192$ com $q < 2^{62}$ atingem 128 bits
de segurança.

\paragraph{A biblioteca \texttt{ckks\_core}.}
O \emph{backend} C++ \texttt{ckks\_core}, implementado para suportar os experimentos
da dissertação, opera em regime somente-aditivo com um único primo $q$ em vez de uma
cadeia RNS completa. Essa escolha simplifica a implementação e reduz o \emph{overhead}
de cifragem ao mínimo necessário para a agregação federada, onde nenhuma multiplicação
entre textos cifrados é realizada. As multiplicações polinomiais internas (cifragem e
decifragem) usam a Transformada Numérica Teórica (NTT) com complexidade $O(N\log N)$,
e toda a aritmética modular é implementada com redução de Montgomery para eliminar
divisões do caminho crítico. Os detalhes de implementação de baixo nível estão
documentados no relatório técnico complementar. A interface Python (\texttt{ckks\_python})
e a camada de alto nível para Flower (\texttt{ckks\_federated}) expõem esses
mecanismos de forma transparente ao praticante de FL.

\paragraph{Síntese.}
O CKKS homomórfico aditivo é uma solução técnica concreta para o problema de privacidade
na agregação federada. Seus três parâmetros essenciais, $N$, $\Delta$ e $K$, devem ser
calibrados em conjunto: $N$ determina a capacidade de empacotamento e o custo de banda;
$\Delta$ controla a precisão; $K$ amplifica o ruído por $\sqrt{K}$. Para as dimensões
típicas de modelos de FL (dezenas de milhares a poucos milhões de parâmetros) e federações
de até algumas centenas de clientes, os parâmetros $N=8192$ e $\Delta=2^{40}$ a $2^{50}$
produzem ruído agregado abaixo de $10^{-4}$ por slot, compatível com a precisão de
treinamento de modelos de aprendizado de máquina.

% =============================================================================
\begin{thebibliography}{9}

\bibitem{Cheon2017}
J.~H. Cheon, A.~Kim, M.~Kim, e Y.~Song,
``Homomorphic encryption for arithmetic of approximate numbers,''
em \emph{ASIACRYPT 2017}, LNCS vol.~10624, Springer, pp.~409--437.

\bibitem{LPR2010}
V.~Lyubashevsky, C.~Peikert, e O.~Regev,
``On ideal lattices and learning with errors over rings,''
em \emph{EUROCRYPT 2010}, LNCS vol.~6110, Springer, pp.~1--23.

\bibitem{sealmanual}
Microsoft Research,
\emph{SEAL: Simple Encrypted Arithmetic Library -- Manual}, v.~4.1, 2023.
\url{https://github.com/microsoft/SEAL}

\bibitem{Geiping2020}
J.~Geiping, H.~Bauermeister, H.~Dröge, e M.~Moeller,
``Inverting gradients -- How easy is it to break privacy in federated learning?''
em \emph{NeurIPS 2020}.

\bibitem{Shokri2017}
R.~Shokri, M.~Stronati, C.~Song, e V.~Shmatikov,
``Membership inference attacks against machine learning models,''
em \emph{IEEE S\&P 2017}.

\end{thebibliography}

\end{document}
